\section{Caso aplicado: asistente de facturación}
\label{sec:caso}
\begin{respuesta}
Escogí el asistente de facturación de servicio al cliente porque concentra lo difícil en un solo
agente: datos confidenciales, conocimiento regulado que cambia cada mes, cifras que no pueden inventarse y
una frontera clara entre resolver y escalar. Lo contrasto con cobranzas, donde la ley
obliga a que el flujo lo decida el código.
\end{respuesta}

Un cliente empresarial pregunta \emph{¿por qué subió mi factura?} La respuesta exige su
consumo medido, su factura, y la explicación de los componentes del costo unitario
(generación, transmisión, distribución, comercialización, pérdidas, restricciones). La empresa tiene
una ventaja: los datos del medidor inteligente explican la variación con precisión.

\begin{figure}[!htbp]
  \centering
  \begin{tikzpicture}[x=1cm,y=1cm]
\node[neut, minimum height=0.75cm] (u) at (0,0) {\textbf{Cliente}\\autenticado};
\node[ctrl, text width=2.0cm, minimum height=0.95cm] (id) at (2.65,0) {\textbf{Identidad}\\\texttt{cliente\_id} de la sesión};
\node[risk, text width=2.0cm, minimum height=0.95cm] (gi) at (5.3,0) {\textbf{Guard de entrada}\\inyección, alcance, longitud};
\node[exec, text width=2.2cm, minimum height=1.2cm] (ag) at (8.3,0) {\textbf{Agente}\\tier \texttt{estandar}\\plantilla asistente-con-datos};
\node[risk, text width=2.2cm, minimum height=0.95cm] (go) at (11.4,0) {\textbf{Guard de salida}\\cifras verificadas, PII ajena, promesas};
\node[neut, minimum height=0.75cm] (r) at (14.3,0) {\textbf{Respuesta}\\citada};
\draw[flecha] (u) -- (id); \draw[flecha] (id) -- (gi); \draw[flecha] (gi) -- (ag); \draw[flecha] (ag) -- (go); \draw[flecha] (go) -- (r);
% herramientas
\node[dat, text width=2.3cm, font=\sffamily\tiny] (t1) at (4.75,-2.3) {\textbf{consultar\_factura}\\lectura};
\node[dat, text width=2.3cm, font=\sffamily\tiny] (t2) at (7.3,-2.3) {\textbf{consultar\_consumo}\\medidor, lectura};
\node[dat, text width=2.3cm, font=\sffamily\tiny] (t3) at (9.85,-2.3) {\textbf{calcular\_desglose}\\determinista};
\node[exec, text width=2.3cm, font=\sffamily\tiny] (t4) at (12.5,-2.3) {\textbf{buscar\_conocimiento}\\tarifas, procedimientos, FAQ};
\foreach \t in {t1,t2,t3,t4}{\draw[flecha, {Stealth[length=4pt]}-{Stealth[length=4pt]}] (ag.south) -- (\t.north);}
\node[obs, text width=2.0cm, font=\sffamily\tiny] (t5) at (1.9,-2.3) {\textbf{crear\_caso}\\escritura: requiere confirmación};
\draw[flecha] (ag.south west) to[out=-150, in=30] (t5.north east);
\node[neut, text width=2.0cm] (h) at (1.9,-3.9) {\textbf{Asesor humano}\\recibe resumen};
\draw[flecha] (t5) -- node[etq, right]{escalar} (h);
\node[plain, text=docgris, font=\sffamily\tiny\itshape] at (8.6,-3.3) {todas las herramientas reciben el identificador del cliente desde el runtime, nunca desde el modelo};
\begin{scope}[on background layer]
\node[grupo, draw=riskbd!60, fit=(id)(gi)(go), inner sep=4pt, label={[font=\sffamily\tiny, text=riskbd]above:capas heredadas de la plataforma}] {};
\end{scope}
\end{tikzpicture}

  \caption{Grafo del asistente de facturación construido desde su manifiesto.}
  \label{fig:facturacion}
\end{figure}

\decision{El sujeto viene de la sesión y las cifras solo de herramientas.}
{``Dame la factura de la empresa X'' no se rechaza con inteligencia: la arquitectura lo hace
imposible. Y el guard de salida verifica que toda cifra de la respuesta aparezca en algún
resultado de herramienta de esa conversación.}
{Confiar en que las instrucciones del modelo eviten estos errores.}

\begin{table}[!htbp]
\centering\small
\caption{Cada pregunta va al mecanismo de su tipo de conocimiento (\cref{fig:conocimiento}).}
\label{tab:preguntas}
\renewcommand{\arraystretch}{1.1}
\begin{tabularx}{\textwidth}{@{}L{4.6cm}L{4.6cm}Y@{}}
\toprule
\textbf{Pregunta} & \textbf{Mecanismo} & \textbf{Por qué}\\
\midrule
¿Por qué subió mi factura? & Consumo + factura + RAG & Datos propios y explicación regulada\\
¿Qué es el componente de distribución? & RAG con vigencia & Conocimiento extenso y citable\\
¿Cuánto debo y cuándo vence? & \texttt{consultar\_factura} & Dato del cliente\\
¿Cuánto pagaría con 10\,\% menos consumo? & \texttt{calcular\_desglose} & Cálculo determinista\\
Quiero un acuerdo de pago & \texttt{crear\_caso} y escalar & Fuera de su autonomía\\
\bottomrule
\end{tabularx}
\end{table}

\textbf{Éxito de negocio.} Conversaciones resueltas sin asesor, tiempo de respuesta y, la más
honesta, tasa de reapertura: si el cliente vuelve a preguntar lo mismo por otro canal, el
agente trasladó trabajo en vez de ahorrarlo.

\subsection*{Comparación con cobranzas}
La Ley 2300 de 2023 \citep{ley2300} solo permite contactar por canales autorizados, prohíbe
usar varios canales en la misma semana o contactar más de una vez al día, y limita los
horarios a lunes a viernes de 7:00 a.m. a 7:00 p.m. y sábados de 8:00 a.m. a 3:00 p.m.

\begin{figure}[!htbp]
  \centering
  \begin{tikzpicture}[x=1cm,y=1cm, cod/.style={ctrl, text width=2.25cm, minimum height=1.15cm}]
\node[dat, text width=1.9cm, minimum height=1.15cm] (c0) at (0,0) {\textbf{Cartera vencida}\\sistema de cobro};
\node[cod] (c1) at (2.85,0) {\textbf{Reglas Ley 2300}\\canal autorizado, horario, frecuencia};
\node[exec, text width=2.25cm, minimum height=1.15cm] (c2) at (5.85,0) {\textbf{LLM redacta}\\tier \texttt{rapido}, dentro de plantilla};
\node[cod] (c3) at (8.85,0) {\textbf{Validador}\\montos y fechas = sistema; tono};
\node[neut, text width=2.1cm, minimum height=1.15cm] (c4) at (11.8,0) {\textbf{Aprobación humana}\\autonomía L4};
\node[cod] (c5) at (14.65,0) {\textbf{Envío y registro}\\por el canal autorizado};
\foreach \a/\b in {c0/c1,c1/c2,c2/c3,c3/c4,c4/c5}{\draw[flecha] (\a) -- (\b);}
\node[obs, text width=2.25cm] (rp) at (2.85,-2.0) {\textbf{Reprogramar}\\fuera de ventana o ya contactado};
\draw[flecha] (c1) -- node[etq, right]{no apto} (rp);
\draw[flecha, draw=riskbd] (c3.south) -- ++(0,-0.55) -| node[etq, pos=0.25, below, text=riskbd]{no coincide: se regenera o se descarta} (c2.south);
\draw[flecha, draw=riskbd] (c4.north) -- ++(0,0.5) -| node[etq, pos=0.25, above, text=riskbd]{rechazo: comentario a la plantilla} (c2.north);
\node[plain, font=\sffamily\scriptsize\bfseries, text=docacento] at (10.2,-1.95) {El código decide el flujo; el modelo solo redacta.};
\node[ctrl, minimum width=0.4cm, minimum height=0.25cm] (l1) at (7.2,-2.6) {};
\node[plain, anchor=west] at (l1.east) {código determinista};
\node[exec, minimum width=0.4cm, minimum height=0.25cm] (l2) at (10.6,-2.6) {};
\node[plain, anchor=west] at (l2.east) {paso con LLM};
\end{tikzpicture}

  \caption{Flujo de cobranza: el modelo ocupa un solo paso acotado y verificado.}
  \label{fig:cobranzas}
\end{figure}

\decision{Cobranza como flujo determinista con aprobación humana (L4).}
{Las reglas legales no pueden depender de que un modelo ``se acuerde'' de ellas: viven en
código y se prueban con casos deterministas. Cada paso con permisos distintos es un componente
separado: quien lee la cartera no puede enviar.}
{Un agente autónomo de cobranza con las reglas en su prompt.}
